\begin{abstract}
% Hard to generalize to open-world and complicated tasks -> abstract

% Online, one-time challenging

% Not flexible, so need progressive distillation and batch evolution

% multiple agentic datasets and settings

% open-end task complexity, generalization, online

% Q1, Q2, Q3, Q4, Q5 
% Q1 -> T1 
% H_0 -> H_1 

% H_1 -> Q2 -> T2 
% H_2 

% H_0 -> Q1 - T1 T1.1 -> T1.2

% zero-start

% Large language model (LLM) agents have shown strong performance in complex task-solving settings, such as web interaction, software engineering, tool use, and long-horizon reasoning. However, most agents still operate as one-off solvers: after executing a task and failing, they often move on without turning the trajectory into reusable lessons for future tasks \tx{可以说不吸取教训重复失误}. This limitation becomes critical \tx{转折有点奇怪} in realistic online settings, where agents face a continuous stream of new tasks rather than repeated trials of the same task. The agent, therefore, needs to learn from each task as it happens. Recent work has explored reflection, memory, rules, and skills as external mechanisms for reusing agent experience. These mechanisms suggest that agent learning does not need to happen only through parameter updates. Instead, an agent can improve by updating an external harness that stores reusable knowledge for future tasks. We formulate this setting as \emph{online harness learning}, where a frozen agent adapts over time by maintaining and updating a structured harness medium \tx{更强调online和真实task一点}. In this work, we instantiate the harness with skills, since skills encode operational knowledge that can directly guide future planning, tool use, verification, and recovery \tx{我们用skill因为可以把各种learning融合}. We propose \emph{context-to-harness skill compilation}, a staged \tx{不要说staged} process that converts one-time execution context into reusable skill harness updates. The solver first proposes candidate skills with the full context. An evolver then abstracts, compares, and updates these proposals against existing skill folders, compiling them into complementary harnesses: general skills for transferable cross-task patterns and topic skills for task-grounded operational knowledge. \tx{说一下架构context solver的reflection逐步compile到evolver的complementary skill harness，并且是from scratch并且和harness interact的} We instantiate the same framework across several realistic benchmark settings including Terminal-Bench, xxxx with sequential task batches, and analyze how evolver designs, feedback types, and harness transfer settings affect online improvement. Our results show that skill harnesses provide a lightweight and interpretable carrier for online adaptation, enabling frozen agents to reuse one-time experience and improve over task streams without parameter updates.

% Our results show that skill harnesses provide a lightweight and interpretable medium for online adaptation, enabling frozen agents to reuse one-time experience and improve over task streams without parameter updates.

% where a frozen agent improves by continually updating a structured harness across sequential tasks. This formulation allows us to systematically investigate how key factors, including evolver designs, feedback types, and harness transfer settings, impact continuous self-improvement. Driven by this objective, we introduce \emph{Evo-Harness}, an instantiation that transforms heterogeneous execution contexts and learning signals into actionable guidance. At its core is \emph{context-to-harness skill compilation}, which incrementally distills useful lessons from noisy, online executions into a reusable skill harness, supporting both broad cross-domain transfer and topic-level adaptation. 

Learning from experience is critical for developing capable, self-improving large language model (LLM) agents. Existing methods typically extract knowledge from accumulated trajectories via reflection, memory, rules, or skills. However, agents in realistic environments continuously encounter novel tasks, often offering only a one-shot opportunity to improve. These executions yield rich but highly noisy contexts, entangling broadly useful lessons with task-specific artifacts. Critically, prior works rarely validate their effectiveness on complex real-world tasks or isolate the underlying drivers of improvement. To address these gaps, we formulate \emph{online harness learning}, where a frozen agent improves by continually updating a structured harness across sequential tasks. This formulation enables a systematic study of key self-improvement factors through our proposed \emph{Evo-Harness}. At its core, \emph{context-to-harness skill compilation} distills noisy, single-shot executions into reusable skill harnesses for cross-domain and topic-level adaptation. To demonstrate the efficacy of one-shot skill compilation, we evaluate across five realistic benchmarks (TerminalBench2, SWE-bench, CL-Bench, $\tau$-bench, WebArena-Infinity). Our extensive analysis demonstrates the effectiveness of \algname{} and provides a principled understanding of how LLM agents can effectively learn on the fly. Our code is available at
\href{https://github.com/A-EVO-Lab/a-evolve/tree/release/evo-harness}{Link}.

% \xy{should we put on the fly in the title, "on the fly context to harness skill compilation ..."}

% Large language model (LLM) agents have shown strong performance in complex task-solving settings, such as web interaction, software engineering, tool use, and long-horizon reasoning. However, when an agent fails on a task, the failure often remains an isolated event rather than a source of future improvement. Many experience-based methods address this issue by learning from accumulated trajectories, where repeated or collected executions help reveal stable patterns. However, in realistic online settings, agents continuously encounter new tasks and contexts, and each completed execution \tingwei{might want to say the limitation of those "revealed patterns" discovered by prior methods instead of the execution itself?} \tingwei{also, the online setting can be put upfront? before we discuss exp-based methods} may provide only a limited opportunity to improve future behavior. These executions produce rich but noisy contexts, where useful lessons are mixed with task-specific details.\tingwei{"task-specific details" sound not very detrimental? e.g. artifacts} Prior work has explored reflection, memory, rules, and skills as media for reusing agent experience, suggesting that adaptation can happen outside model parameters.
% \tingwei{(i) are these the same set of prior work as the exp-based method above? (ii) if not, should we mention the limitation of these inference-only methods as well?}
% Building on this view, we formulate this setting as \emph{online harness learning}, where a frozen agent improves by incrementally \tingwei{iteratively/continually?} updating a structured harness that is reused in later tasks. We instantiate this harness with skills to integrate diverse context executions and learning signals from realistic agent tasks into actionable guidance. We propose \emph{context-to-harness skill compilation}, which incrementally compiles lessons from individual execution contexts into reusable skill harness updates for both broad transfer and topic-level adaptation. We instantiate the framework across several realistic benchmark settings, including Terminal-Bench, SWE-Bench, Tau-Bench, CL-Bench, and related settings, using sequential task batches. We further analyze how evolver designs, feedback types, and harness transfer settings affect online improvement. 
% \hq{suggested another abstract framing: methodology improvement is marginal, but several counter-intuitive findings are rare for self-evolving agents development. Self-evolution gain is not free—— It depends on feedback grounding、skill abstraction and solver-evolver alignment，in the worst case, online evolution is worse than evolution。EvoHarness is a pipeline that developed on these three findings and empirical supports。}

% \hq{Large language model (LLM) agents are increasingly evaluated in online settings where each task execution is a one-shot opportunity for improvement. A growing line of work treats this as a method problem — proposing reflection, memory, rules, and skill pipelines to convert executions into reusable guidance. In this work, we argue that the method-level differences between these pipelines are not the main driver of online self-evolution. Through a controlled study across five benchmarks (CL-Bench, TerminalBench-2, SWE-bench Lite, τ-bench, WebArena-Infinity) and five foundation models, we find that three under-examined factors largely determine whether online evolution helps at all: (i) feedback grounding — self-generated feedback can drop performance below the no-evolution baseline; (ii) skill abstraction level — broad versus topic-localized guidance contribute asymmetrically across task types; (iii) solver-evolver alignment — cross-model skill transfer is sensitive to reasoning-pattern compatibility and can be net-negative. To enable this study we build Evo-Harness, a propose-curate pipeline that compiles execution contexts into a structured skill harness; configured according to our findings, Evo-Harness improves five frozen agents over six prior skill/memory baselines on every benchmark (+2.3 to +6.7 over the strongest baseline). We release the pipeline as an analytical instrument for studying when, why, and for whom online self-evolution works.}
\end{abstract}
